\SourcePage{193}{196}
\section{Transitions within the continuous spectrum}

One of the most important uses of perturbation theory is the calculation of
transition probabilities within the continuous spectrum under a constant
(time-independent) perturbation. As already noted, continuous-spectrum states
are almost always degenerate. After choosing a definite set of unperturbed
wave functions for a given energy level, the question can be posed as follows.
The system is known to have occupied one of these states initially; we seek
the probability that it passes into another state of the same energy. For
transitions from an initial state $i$ into states between $\nu_f$ and
$\nu_f+d\nu_f$, (42.5), with $\omega=0$ and a change of notation, gives
directly
\begin{equation}
  dw_{fi}=\frac{2\pi}{\hbar}|V_{fi}|^2
  \delta(E_f-E_i)\,d\nu_f.
  \tag{43.1}
\end{equation}
As expected, this is nonzero only when $E_f=E_i$: a constant perturbation
causes transitions solely between states of equal energy. For transitions
from continuous-spectrum states, however, $dw_{fi}$ cannot itself be treated
as a transition probability. It does not even have the corresponding
dimension, but has dimension $1/\mathrm{s}$. Expression (43.1) gives the
number of transitions per unit time, and its dimension depends on the chosen
normalization of the continuous-spectrum wave functions.\footnote{Phenomena
covered by this theory include, for example, various collisions. In their
initial and final states the system is an assemblage of free particles, while
their mutual interaction serves as the perturbation. With a suitable wave-
function normalization, (43.1) can then become a collision cross-section
(see \secref{126}).}

We shall calculate the perturbed wave function which, before the perturbation
begins to act, coincides with the original unperturbed
\SourcePage{194}{197}
function $\psi_i^{(0)}$. Following the procedure stated at the end of the
preceding section, take the perturbation to be switched on adiabatically as
$e^{\lambda t}$, with $\lambda\to0$. Formula (42.7), after setting
$\omega=0$ and changing notation, gives
\begin{equation}
  a_{fi}^{(1)}=V_{fi}
  \frac{\exp\left\{\dfrac{i}{\hbar}(E_f-E_i)t+\lambda t\right\}}
       {E_i-E_f+i\lambda}.
  \tag{43.2}
\end{equation}
The perturbed wave function is
\[
  \Psi_i=\Psi_i^{(0)}+\int a_{fi}^{(1)}\Psi_f^{(0)}\,d\nu_f,
\]
where the integration extends over the entire continuous
spectrum.\footnote{If a discrete spectrum is also present, the corresponding
sum over its states must be added to the integral here and in the formulas
below.} Substitution of (43.2) yields
\begin{equation}
  \Psi_i=\left[\psi_i^{(0)}+\int
  V_{fi}\psi_f^{(0)}\frac{d\nu_f}{E_i-E_f+i0}\right]
  \exp\left(-\frac{i}{\hbar}E_it\right).
  \tag{43.3}
\end{equation}
In the limit $\lambda\to0$, the factor $e^{\lambda t}$ has been replaced by
unity. The term $+i0$, meaning the limit of $i\lambda$ as the positive
quantity $\lambda$ tends to zero, specifies the integration prescription for
$E_f$. Its differential occurs as a factor in $d\nu_f$, together with those
of the other quantities characterizing continuous-spectrum states. Without
$i\lambda$, the integrand in (43.3) would have a pole at $E_f=E_i$, and the
integral would diverge near it. The term $i\lambda$ displaces this pole into
the upper half-plane of the complex variable $E_f$. When $\lambda\to0$, the
pole returns to the real axis, but the integration path is now known to pass
below it:
\begin{equation}
\begin{tikzpicture}[x=1cm,y=1cm,>=stealth,baseline=-0.35ex]
  \draw[->] (-2.8,0) -- (-1.25,0);
  \draw (-1.25,0) -- (-0.48,0)
        arc[start angle=180,end angle=360,radius=0.48]
        -- (1.25,0);
  \draw[->] (1.25,0) -- (2.8,0);
  \fill (0,0) circle (1.5pt);
  \node[above=2pt] at (0,0) {$E_i$};
  \node[below=2pt] at (2.45,0) {$E_f$};
\end{tikzpicture}
  \tag{43.4}
\end{equation}
The time factor in (43.3) shows, as it should, that this function belongs to
the same energy $E_i$ as the original unperturbed function. Equivalently,
\begin{equation}
  \psi_i=\psi_i^{(0)}+\int
  \frac{V_{fi}}{E_i-E_f+i0}\psi_f^{(0)}\,d\nu_f
  \tag{43.5}
\end{equation}
\SourcePage{195}{198}
satisfies the Schrödinger equation
\[
  (\hat H_0+\hat V)\psi_i=E_i\psi_i.
\]
It is consequently natural that this expression agrees exactly with (38.8).\footnote{The integration prescription in (43.5) can also be fixed by
requiring the large-distance asymptotic form of $\psi_i$ to contain only an
outgoing wave, and no incoming wave (see \secref{136}).}

The preceding calculation is the first approximation of perturbation theory.
The second approximation is also readily obtained. One derives the next
approximation for $\Psi_i$ by the method of \secref{38}, now knowing how the
``divergent'' integrals are to be taken. A direct calculation gives
\[
  \Psi_i=\left\{\psi_i^{(0)}+\int\left[
  V_{fi}+\int\frac{V_{f\nu}V_{\nu i}}{E_i-E_\nu+i0}\,d\nu
  \right]\frac{\psi_f^{(0)}\,d\nu_f}{E_i-E_f+i0}\right\}
  \exp\left(-\frac{i}{\hbar}E_it\right).
\]

Comparison with (43.3) lets us write the corresponding probability formula
(more precisely, the number of transitions) by direct analogy with (43.1):
\begin{equation}
  dw_{fi}=\frac{2\pi}{\hbar}
  \left|V_{fi}+\int\frac{V_{f\nu}V_{\nu i}}
  {E_i-E_\nu+i0}\,d\nu\right|^2
  \delta(E_i-E_f)\,d\nu_f.
  \tag{43.6}
\end{equation}
The matrix element $V_{fi}$ may vanish for the transition in question. The
first-order effect is then absent altogether, and (43.6) becomes
\begin{equation}
  dw_{fi}=\frac{2\pi}{\hbar}
  \left|\int\frac{V_{f\nu}V_{\nu i}}{E_i-E_\nu}\,d\nu\right|^2
  \delta(E_f-E_i)\,d\nu_f.
  \tag{43.7}
\end{equation}
In applications of this formula, $E_\nu=E_i$ is usually not a pole of the
integrand. The prescription for integration over $dE_\nu$ is then immaterial,
and the path may be taken directly along the real axis.

States $\nu$ for which $V_{f\nu}$ and $V_{\nu i}$ are nonzero are often
called \emph{intermediate} states for the transition $i\to f$. In a pictorial
description, the transition appears to take place in two stages,
$i\to\nu$ and $\nu\to f$; this description must not, of course, be understood
literally. A transition $i\to f$ may require not one but several successive
intermediate states. Formula (43.7)
\SourcePage{196}{199}
extends directly to such cases. If two intermediate states are needed, then
\begin{equation}
  dw_{fi}=\frac{2\pi}{\hbar}
  \left|\int\frac{V_{f\nu'}V_{\nu'\nu}V_{\nu i}}
  {(E_i-E_{\nu'})(E_i-E_\nu)}\,d\nu\,d\nu'\right|^2
  \delta(E_f-E_i)\,d\nu_f.
  \tag{43.8}
\end{equation}

Finally, the mathematical meaning of integrals taken along a path such as
(43.4) is expressed by
\begin{equation}
  \int\frac{f(x)\,dx}{x-a-i0}
  =\operatorname{P}\!\int\frac{f(x)\,dx}{x-a}+i\pi f(a),
  \tag{43.9}
\end{equation}
where the interval of integration on the real axis contains $x=a$. Indeed,
detouring around the pole $x=a$ on a semicircle of radius $\rho$, the full
integral is the sum of the real-axis integrals from the lower limit to
$a-\rho$ and from $a+\rho$ to the upper limit, plus $i\pi$ times the residue
at the pole. As $\rho\to0$, the real-axis pieces combine into the integral
over the full interval understood as a principal value, as indicated by the
marked integral sign, and (43.9) follows. The same formula is also written
symbolically as
\begin{equation}
  \frac1{x-a-i0}=P\frac1{x-a}+i\pi\delta(x-a);
  \tag{43.10}
\end{equation}
the symbol $P$ means that the principal value is to be taken when
$f(x)/(x-a)$ is integrated.
